\begin{mistake}{2026-07-20}{06}

\begin{problemcontent}
设 \(f(x)\) 有连续的二阶导数，曲线 \(y=f(x)\) 在点 \((0,0)\) 处与 \(x\) 轴相切，且 \(\lim_{x\to 0}\frac{f(x)}{x^2}=1\)，求曲线 \(y=f(x)\) 在点 \((0,0)\) 处的曲率圆方程。
\end{problemcontent}

\begin{solutioncontent}
由曲线在 \((0,0)\) 处与 \(x\) 轴相切可知，曲线过原点且该点切线斜率为 \(0\)，即：
\[ f(0) = 0, \quad f'(0) = 0 \]
又因 \(f(x)\) 有连续的二阶导数，利用洛必达法则计算极限：
\[ \lim_{x\to 0}\frac{f(x)}{x^2} = \lim_{x\to 0}\frac{f'(x)}{2x} = \lim_{x\to 0}\frac{f''(x)}{2} = \frac{f''(0)}{2} \]
已知该极限为 \(1\)，故 \(f''(0) = 2\)。
在点 \((0,0)\) 处，曲率半径为：
\[ R = \frac{(1+[f'(0)]^2)^{\frac{3}{2}}}{|f''(0)|} = \frac{1}{2} \]
因 \(f'(0)=0\)，法线为 \(y\) 轴；又因 \(f''(0) = 2 > 0\)，曲线在原点处凹（开口向上），故曲率圆圆心位于 \(y\) 轴正半轴上，距离原点为 \(R = \frac{1}{2}\)。
圆心坐标为 \((0, \frac{1}{2})\)，曲率圆方程为：
\[ x^2 + \left(y - \frac{1}{2}\right)^2 = \frac{1}{4} \]
\end{solutioncontent}

\mistakeknowledge{
\begin{itemize}
  \item \textbf{核心公式}：曲率半径 \(R = \frac{(1+(y')^2)^{\frac{3}{2}}}{|y''|}\)。
  \item \textbf{易错点}：曲率圆圆心位置的确定。当 \(y'=0\) 时无需硬背圆心坐标公式，结合凹凸性（\(y''>0\) 向上偏）与法线方向直接推导更高效。
  \item \textbf{关联知识}：曲线与 \(x\) 轴相切的代数等价条件是 \(f(x_0)=0\) 且 \(f'(x_0)=0\)。
\end{itemize}}

\end{mistake}
